% Engine-owned annotation primitives. Inlined into each helper-using document.
\usetikzlibrary{decorations.markings,shapes.arrows}

% Hollow block arrow printed between figure panels (the exam glyph U+21E8 is missing from the TeX fonts).
% #2 is a coordinate such as 4.7,1 or a named point.
\newcommand{\ExamImplies}[2][]{%
  \node[single arrow,draw,line width=.5pt,minimum height=7mm,inner sep=.8mm,
        single arrow head extend=.9mm,single arrow tip angle=90,#1] at (#2) {};%
}

% B is the corner; the caller must confirm a right-angle mark in the source.
\newcommand{\ExamRightAngle}[4][]{%
  \pic[draw,angle radius=2mm,#1] {right angle=#2--#3--#4};%
}

% Side of A->B on which point #3 lies: \ExamSide is 1 on the left, -1 on the right (or on the line).
\newcommand{\ExamSideOf}[3]{%
  \path let \p1=(#1),\p2=(#2),\p3=(#3) in
    \pgfextra{\pgfmathparse{((\x2-\x1)/100*(\y3-\y1)/100-(\y2-\y1)/100*(\x3-\x1)/100)>0.0001 ? 1 : -1}%
      \xdef\ExamSide{\pgfmathresult}};%
}

% Side of A->B on which the figure lies, as \ExamSide (1 left, -1 right): the centre of the picture's named
% points defined so far (\ExamPoints, listed by the engine). The bounding box alone cannot tell: its centre
% lies on the hypotenuse of every right triangle. It is used only when fewer than three points are named or
% their centre lies on the line AB.
\providecommand{\ExamPoints}{}
\makeatletter
\newcommand{\ExamInnerSide}[2]{%
  \gdef\exam@sx{0}\gdef\exam@sy{0}\gdef\exam@count{0}%
  \@for\exam@p:=\ExamPoints\do{\pgfutil@ifundefined{pgf@sh@ns@\exam@p}{}{%
    \path let \p1=(\exam@p) in \pgfextra{%
      \pgfmathparse{\exam@sx+\x1/100}\xdef\exam@sx{\pgfmathresult}%
      \pgfmathparse{\exam@sy+\y1/100}\xdef\exam@sy{\pgfmathresult}%
      \pgfmathparse{\exam@count+1}\xdef\exam@count{\pgfmathresult}};}}%
  \ExamSideOf{#1}{#2}{current bounding box.center}%
  \ifdim\exam@count pt>2.5pt
    \path let \p1=(#1),\p2=(#2) in \pgfextra{%
      \pgfmathparse{(\x2-\x1)/100*(\exam@sy/\exam@count-\y1/100)-(\y2-\y1)/100*(\exam@sx/\exam@count-\x1/100)}%
      \xdef\exam@cross{\pgfmathresult}%
      \pgfmathparse{abs(\exam@cross)>0.03*(((\x2-\x1)/100)^2+((\y2-\y1)/100)^2) ? (\exam@cross>0 ? 1 : -1) : \ExamSide}%
      \xdef\ExamSide{\pgfmathresult}};%
  \fi
}

% Dashed length curve from A to B, bowed away from the figure (\ExamInnerSide), so the order of A and B does
% not matter. [inside] bows into the figure; an explicit bend left/bend right is kept.
% No arrowheads unless explicitly requested.
\tikzset{inside/.style={}}
\newcommand{\ExamLengthArc}[4][]{%
  \ExamInnerSide{#2}{#3}%
  \pgfutil@in@{inside}{#1}%
  \ifpgfutil@in@\pgfmathparse{-\ExamSide}\xdef\ExamSide{\pgfmathresult}\fi
  \ifdim\ExamSide pt>0pt\def\ExamArcBend{bend right}\else\def\ExamArcBend{bend left}\fi
  % Printed length arcs rise 0.12 of the length they span (25 source arcs read on the scans: quartiles 0.10,
  % 0.12, 0.13; the fixed 18 degree bend rose 0.08), and a short one needs room for its label between arc
  % and line: rise = max(0.12 of the span, 2.2mm), at most a 45 degree bend.
  \path let \p1=($(#3)-(#2)$), \n1={veclen(\x1,\y1)} in
    \pgfextra{\pgfmathparse{min(45,2*atan(2*max(0.12*\n1,6.3)/max(\n1,1)))}\xdef\ExamArcAngle{\pgfmathresult}};%
  \draw[dashed,\ExamArcBend=\ExamArcAngle,#1] (#2) to
    node[midway,fill=white,inner sep=.6pt] {#4} (#3);%
}
\makeatother

% Angle mark at vertex B between BA and BC, always the smaller angle whatever the order of A and C.
% #5 is the printed label ({} for none); [double] doubles the arc, [angle radius=5mm] resizes it,
% [angle eccentricity=.6] puts a mark such as $\bullet$ or $\times$ inside the arc.
\newcommand{\ExamAngle}[5][]{%
  \ExamSideOf{#3}{#2}{#4}%
  % A narrow angle has no room for its label near the vertex: the label moves out along the angle.
  \path let \p1=(#2),\p2=(#3),\p3=(#4) in
    \pgfextra{\pgfmathparse{abs(mod(atan2(\y1-\y2,\x1-\x2)-atan2(\y3-\y2,\x3-\x2)+540,360)-180)}%
      \pgfmathparse{min(4,max(1.7,0.5/tan(max(\pgfmathresult,2)/2)+0.3))}\xdef\ExamEccentricity{\pgfmathresult}};%
  \ifdim\ExamSide pt>0pt
    \pic[draw,angle radius=3.5mm,angle eccentricity=\ExamEccentricity,inner sep=1pt,"#5",#1] {angle=#2--#3--#4};%
  \else
    \pic[draw,angle radius=3.5mm,angle eccentricity=\ExamEccentricity,inner sep=1pt,"#5",#1] {angle=#4--#3--#2};%
  \fi
}

% Explicit placement, not an automatic collision solver. Never hides lines.
\newcommand{\ExamLabel}[3][above right]{%
  \node[inner sep=2pt,outer sep=0pt,#1] at (#2) {#3};%
}

% Midpoint ticks in the local direction of AB; does not draw the segment.
\newcommand{\ExamTicks}[4][]{%
  \begingroup
  \ifnum#4<1\relax
    \PackageError{exam-marks}{ExamTicks count must be 1, 2 or 3}{Use one of these counts.}%
  \else\ifnum#4>3\relax
    \PackageError{exam-marks}{ExamTicks count must be 1, 2 or 3}{Use one of these counts.}%
  \else
    \path[postaction={decorate},decoration={markings,
      mark=at position .5 with {%
        \foreach \ExamTickIndex in {1,...,#4}{%
          \draw[line width=.4pt,#1]
            ({(\ExamTickIndex-(#4+1)/2)*2.2pt},-1.8pt) --
            ({(\ExamTickIndex-(#4+1)/2)*2.2pt},1.8pt);%
        }%
      }}] (#2)--(#3);%
  \fi\fi
  \endgroup
}
